\documentclass[article, 11pt]{article}

%%%%%%%%%%%%%%%%%%%%%%%%%%%%%%%%%%%%%%%%%%%%%%%%%%%
% Preamble:
%%%%%%%%%%%%%%%%%%%%%%%%%%%%%%%%%%%%%%%%%%%%%%%%%%%
% Typical Packages:
\usepackage{fullpage}
\usepackage{amsfonts}
\usepackage{amsmath}
\usepackage{amsthm}
\usepackage{graphicx}
\usepackage{color}
\usepackage{palatino, url, multicol}
\usepackage{enumerate}
\usepackage{ulem}
\usepackage{hyperref}
\usepackage{verbatim}

\thispagestyle{empty}

%%%%%%%%%%%%%%%%%%%%%%%%%%%%%%%%%%%%%%%%%%%%%%%%%%%%
% Start Document:
%%%%%%%%%%%%%%%%%%%%%%%%%%%%%%%%%%%%%%%%%%%%%%%%%%%%
\begin{document}

% Project Title:
\begin{center}
\boxed{\LARGE{\textbf{\textsc{MINI-PROJECT 1}}}}\\
\Huge{\textbf{\textsc{Optimization Using Calculus}}}
\end{center}

\vspace{.6cm}

% Brief Summary:
\noindent \boxed{\textbf{\textsc{Brief Summary}}}

\vspace{-.35cm}
\noindent \hrulefill

Optimization is the process of finding the best solution from a set of feasible solutions to a problem, typically by maximizing or minimizing a specific objective function while satisfying certain constraints. The goal of this mini-project is to streamline this basic exercise in calculus and to expand on it in several ways by introducing one or more parameters. Formulate a problem of your choice and optimize the objective function. Once the optimal variables are found, analyze how the model variables or model outputs depend on parameters in the system. Present your findings and results in a written report. \\

% Goals:
\noindent \boxed{\textbf{\textsc{Goals}}}

\vspace{-.35cm}
\noindent \hrulefill

\begin{itemize}
	\item Formulate and solve (by hand) the general optimization problem for fixed values of the parameters.
	\item Write Python code that solves the problem, investigates how the solution depends on parameters, and produces figures that graphically display these dependencies.
	\item Provide a written \LaTeX\, report that formally presents your problem and solution.
\end{itemize}

% LaTeX Deliverables:
\noindent \boxed{\textbf{\textsc{\LaTeX\, Deliverables:}}}

\vspace{-.26cm}
\noindent \hrulefill

\begin{itemize}
	\item The report structure follows the provided template.
	\item Include at least one figure (with a caption) that shows how your optimal solution depends on system parameters.
	\item Provide the methodology used to solve the problem, showing the steps that lead to the optimal solution.
\end{itemize}

% Python Deliverables:
\noindent \boxed{\textbf{\textsc{Python Deliverables:}}}

\vspace{-.35cm}
\noindent \hrulefill

\begin{itemize}
	\item A Python function that takes system variables/parameters as input and outputs the optimal solution.
	\item Include at least one well-designed plot with an appropriate title, axis labels, and a legend (as needed).
	\item You must use a built-in Python function such as \texttt{optimize.minimize}.
\end{itemize}

% AI Disclosure:
\noindent \boxed{\textbf{\textsc{AI Disclosure:}}}

\vspace{-.35cm}
\noindent \hrulefill

If you use AI in any way, you must include an Appendix in the write-up that notes the LLM used and provides a list of the prompts you used. You may use AI for suggestions related to Python code, \LaTeX\, syntax, or other general ideas, but \textit{the report should not be copy-and-pasted from AI.}

\end{document}